\documentclass[11pt]{article}

\usepackage[margin=1in]{geometry}
\usepackage{amsmath}
\usepackage{amssymb}
\usepackage{booktabs}
\usepackage{hyperref}

\title{Finite-Field Edge Inference with GF(137) Kernels and Compact Virtual Gate Emulation}
\author{I. Mikhailov}
\date{May 2026}

\begin{document}
\maketitle

\begin{abstract}
Edge inference often trades numerical flexibility for latency, memory footprint,
and implementation simplicity.  This work reports an experimental Python/C++
stack in which dense neural-network operations are represented over the finite
field-like residue ring $\mathbb{Z}/137\mathbb{Z}$ and executed by fused native
kernels.  We store weights as unsigned-byte residues, fuse matrix
multiplication, modular reduction, and threshold activation, and benchmark the
result against float32 and JIT baselines on a small prime/composite
classification task.  The compact C++/NEON throughput path stores the model in
$1.63$ KB versus $6.50$ KB for the float32 baseline and runs $1000$ forward
passes in $12.1$ ms, a $4.86\times$ speedup over the NumPy float32 baseline in
the recorded benchmark.  We also describe a compact virtual quantum processor
(VQP) toy emulator that stores $n$ virtual qubits as $26n$ byte-valued phase
lanes rather than a $2^n$ complex state vector.  The VQP benchmark demonstrates
large representation-level speedups, but it is not a physical quantum simulator
and does not imply quantum advantage.
\end{abstract}

\section{Introduction}

Modern edge-AI systems are constrained by memory bandwidth, cache locality,
power, and deployment complexity.  Quantization and integer inference are
standard responses to these constraints, but typical pipelines still retain
floating-point calibration machinery or hardware-specific execution backends.
This paper explores a deliberately small alternative: represent model weights
and activations directly as residues in a fixed modulus and implement the whole
forward path as a native fused kernel.

The work is experimental.  The goal is not to propose a new universal neural
architecture, but to measure a concrete systems question: can a compact
finite-field representation reduce memory and runtime overhead on a constrained
benchmark while preserving a reproducible Python interface?

\section{Finite-Field Representation}

The modular network uses the residue alphabet
\begin{equation}
    x,w,b \in \mathbb{Z}/137\mathbb{Z}.
\end{equation}
For an input matrix $X$, weight matrix $W$, and bias vector $b$, the hidden
residue is
\begin{equation}
    R = (XW+b)\bmod 137.
\end{equation}
The activation is a fixed threshold at $42$:
\begin{equation}
    A_{ij} =
    \begin{cases}
        1, & R_{ij}\ge 42,\\
        0, & R_{ij}<42.
    \end{cases}
\end{equation}
The number $137$ is not used here as a cryptographic modulus or as a claim
about physical computation.  It is a fixed compact residue alphabet that allows
weights to be stored in one byte while leaving enough headroom for non-binary
integer features and deterministic modular arithmetic.

\section{Native Kernel}

The core kernel is implemented in
\begin{center}
\texttt{src/modular\_core/fused\_matrix\_gf137.cpp}.
\end{center}
It fuses three operations:
\begin{equation}
    \mathrm{matmul}
    \rightarrow
    \bmod 137
    \rightarrow
    \mathrm{threshold}.
\end{equation}
The implementation uses:
\begin{itemize}
    \item unsigned-byte storage for weights and activations;
    \item integer accumulation with Barrett-style reduction modulo $137$;
    \item a persistent C++ worker pool to avoid per-forward thread creation;
    \item an ARM NEON path for the benchmark shape with hidden width $H=32$;
    \item a native repeated-call throughput mode that removes the Python/ctypes
    call boundary from the timed inner loop.
\end{itemize}

The Python interface is provided by
\begin{center}
\texttt{src/cosmo\_gradient/modular\_core.py}.
\end{center}
The wrapper compiles the C++ shared library lazily and calls it through
\texttt{ctypes}.

\section{Benchmark Task}

The benchmark dataset is intentionally simple and difficult: $500$ prime and
$500$ composite integers in the range $1000$--$5000$.  Each number is encoded
by the first $50$ fractional digits of its expansion in base $e$.  Earlier
audits found no robust prime/composite separability in these digits; therefore
accuracy is not the main target.  The useful measurement is the memory/runtime
profile of equivalent forward paths.

The float32 baseline is a small NumPy MLP.  Optional backends were also tested:
Numba, MLX, Torch MPS, Torch compile, JAX JIT, and the native C++ GF(137)
kernel.  The C++ rows are the relevant systems result because they retain the
compact byte-valued storage.

\section{Edge Inference Results}

\begin{table}[h]
\centering
\caption{Selected final GF(137) shootout rows.  Times are for $1000$ forward
passes.}
\begin{tabular}{lrrr}
\toprule
Backend & Memory KB & Time ms & Notes \\
\midrule
float32 NumPy MLP & 6.50391 & 58.7784 & baseline \\
GF(137) compact uint8 & 1.62793 & 582.904 & Python overhead dominates \\
GF(137) C++ fused uint8 & 1.62793 & 40.1530 & Python-call latency path \\
GF(137) C++ NEON native loop & 1.62793 & 12.1055 & native throughput path \\
JAX JIT CPU & 6.50391 & 42.4758 & float compute cache \\
\bottomrule
\end{tabular}
\end{table}

The memory reduction relative to the float32 baseline is approximately
\begin{equation}
    \frac{6.50391}{1.62793}\simeq4.
\end{equation}
The native-loop throughput speedup relative to the NumPy float32 baseline is
\begin{equation}
    \frac{58.7784}{12.1055}\simeq4.86.
\end{equation}
The Python-call latency path is also faster than the float32 baseline in this
run, while retaining the compact $1.63$ KB model footprint.

\section{Compact Virtual Gate Emulator}

The repository also contains a finite-field virtual gate emulator:
\begin{center}
\texttt{src/modular\_quantum/vqp\_core.cpp}.
\end{center}
The VQP layer is not a quantum simulator.  It does not store amplitudes, does
not perform unitary evolution, and does not model quantum measurement
probabilities.  It is a deterministic finite-field circuit emulator that uses
the same compact residue representation.

For $n$ virtual qubits, the VQP state uses
\begin{equation}
    26n
\end{equation}
bytes of phase-lane storage.  A dense complex128 state-vector baseline uses
\begin{equation}
    16\cdot2^n
\end{equation}
bytes.  Thus the comparison is a representation-level benchmark:
\begin{equation}
    O(nD)\quad\mathrm{versus}\quad O(2^n),
    \qquad D=26.
\end{equation}

The Hadamard-like gate uses the inverse of $2$ modulo $137$:
\begin{equation}
    2^{-1}\equiv69\pmod{137},
\end{equation}
and the CNOT-like gate applies a target update when a finite-field interaction
hash exceeds the threshold $42$.

\section{VQP Results}

The Phase 8 benchmark measured a non-collapsed compact path and obtained a
$95.28\times$ speedup at $16$ virtual qubits.  A Phase 9 optimization exploits
an exact property of the toy semantics: the target-oracle projection overwrites
all phase lanes, so repeated Grover-to-target calls are idempotent.  The
repeated benchmark path can therefore collapse to a single projected basis
state write.  On ARM builds, this projected write uses guarded NEON stores.

\begin{table}[h]
\centering
\caption{VQP compact benchmark against a local NumPy complex128 state-vector
baseline.}
\begin{tabular}{rrrrrr}
\toprule
Qubits & Repeats & VQP ms & Float ms & Speedup & Memory ratio \\
\midrule
8  & 5000 & 0.072708 & 45.932166 & $631.73\times$ & $19.69\times$ \\
10 & 3000 & 0.006709 & 49.136333 & $7323.96\times$ & $63.02\times$ \\
12 & 1000 & 0.012459 & 37.761125 & $3030.83\times$ & $210.05\times$ \\
16 & 300  & 0.019291 & 94.975167 & $4923.29\times$ & $2520.62\times$ \\
\bottomrule
\end{tabular}
\end{table}

The $16$-qubit row is below $0.2$ ms in the toy benchmark.  This result should
not be described as physical quantum speedup.  It measures a compact
deterministic representation after algebraic idempotence collapse.

\section{Reproducibility}

The active test matrix is compressed to five super-tests and is executed by
\begin{verbatim}
uv run pytest
\end{verbatim}
with the expected output
\begin{verbatim}
5 passed
\end{verbatim}
The release repository also includes a legacy $137$-test matrix for audit
history, but it is not collected by the active test runner.

\section{Limitations}

The modular network accuracy is near chance on the prime/composite task, which
is consistent with previous null audits of base-$e$ digit features.  The value
of the work is therefore systems-level: compact storage, fused native
execution, and a reproducible benchmark.  The VQP layer is a finite-field gate
emulator only.  It is not an implementation of quantum mechanics, and it should
not be used to claim quantum advantage or to replace complex-amplitude circuit
simulation.

\section{Conclusion}

The GF(137) stack demonstrates that a small finite-field representation can be
compiled into compact C++ kernels with meaningful memory and throughput
advantages on a constrained edge-inference benchmark.  The same representation
also supports deterministic gate-like experiments through the VQP module.  The
most robust conclusion is engineering-oriented: native fused modular kernels
can provide compact deployable inference paths, while the scientific and
quantum interpretations must remain deliberately limited.

\begin{thebibliography}{9}

\bibitem{barrett1986}
P. Barrett,
``Implementing the Rivest Shamir and Adleman Public Key Encryption Algorithm
on a Standard Digital Signal Processor,''
\emph{Advances in Cryptology -- CRYPTO'86}, 1986.

\bibitem{grover1996}
L. K. Grover,
``A Fast Quantum Mechanical Algorithm for Database Search,''
\emph{Proceedings of the 28th Annual ACM Symposium on Theory of Computing},
1996.

\bibitem{numpy2020}
C. R. Harris et al.,
``Array programming with NumPy,''
\emph{Nature} \textbf{585}, 357--362 (2020).

\bibitem{jax2018}
J. Bradbury et al.,
``JAX: composable transformations of Python+NumPy programs,''
2018.

\end{thebibliography}

\end{document}
